\par\SourceLeaf{305}
\indent Überhaupt verhält es sich mit allen solchen neuen Kalkülen so, dass man durch sie
nur leisten kann, was auch ohne sie zu leisten wäre. Der Vorteil ist aber,
dass, wenn ein solcher Kalkül dem innersten Wesen häufig wiederkehrender
Bedürfnisse entspricht, jeder, der ihn sich ganz angeeignet hat, auch ohne die
gleichsam unbewussten Eingebungen des Genies, die niemand \textls{erzwingen} kann, die betreffenden
Aufgaben lösen kann, und zwar selbst in so verwickelten Fällen gleichsam mechanisch,
in denen ohne eine solche Hilfe auch das Genie ohnmächtig wird. So ist es mit
der Erfindung der Buchstabenrechnung überhaupt, so war es mit der Differenzialrechnung,
und so ist es auch (wenn auch in begrenzteren Bereichen) mit \textsc{Lagranges}
Variationsrechnung, mit meiner Kongruenzenrechnung und mit \textsc{Möbius}' Kalkül. Es
werden durch solche Konzeptionen unzählige Aufgaben, die sonst vereinzelt stehen und
jedes Mal neue Anstrengungen (kleinere oder größere) des Erfindungsgeistes erfordern,
gleichsam zu einem organischen Reich. \mbox{.\,.\,.\,.\,.\,.\,.\,.\,.}

\Needspace{9\baselineskip}\begin{center}
\medskip

[II.]

\textbf{[DER BARYZENTRISCHE KALKÜL UND DER RESULTANTENKALKÜL.]}

[1.]
\end{center}

Der baryzentrische Kalkül findet sein Gegenstück in einem anderen (vermutlich noch
umfassenderen) Kalkül, den man Resultantenkalkül nennen könnte. Der erste befasst
sich mit Punkten, in denen man schwere Massen voraussetzt. Der zweite würde
Linien zum Gegenstand haben, auf denen Kräfte wirken. Sind $a$, $b$, $c$,
$d$ usw. solche Linien, auf denen jeweils in einer bestimmten Richtung Kräfte wirken, die
den Zahlen $\alpha$, $\beta$, $\gamma$, $\delta$ usw. proportional sind, so würde
die Gleichung

\[
\alpha a + \beta b + \gamma c + \delta d + \text{usw.} = 0
\]
\noindent bedeuten, dass sich diese Kräfte gegenseitig im Gleichgewicht halten.
